\chapter{線形代数の準備}
\label{ch:found-linalg}

第2章（Gaussian 測度の $W_2$）で用いる線形代数の道具をまとめる．
ベクトルは $\R^n$ の列ベクトルとし，
内積を $\inner{u}{v}:=u^\top v$，ノルムを $\norm{u}:=\sqrt{\inner{u}{u}}$ と書く．
行列・転置 $A^\top$・逆行列・固有値・固有ベクトル・階数 $\rank A$ などの
基本概念は既知とする．
本章で証明なしに認めるのはスペクトル分解
（定理~\ref{thm:la-spectral}）のみで，それ以外の主張には証明を付ける．

\section{対称行列と半正定値性}

\begin{definition}{対称行列・直交行列}{la-symmetric}
  $A\in\R^{n\times n}$ が\textbf{対称}であるとは $A^\top=A$ をみたすことをいう．
  $U\in\R^{n\times n}$ が\textbf{直交行列}であるとは
  $U^\top U=UU^\top=I$（$I$ は単位行列）をみたすことをいう．
  直交行列はノルムを保つ：
  $\norm{Ux}^2=x^\top U^\top Ux=\norm{x}^2$．
\end{definition}

\begin{claim}{正規直交基底への延長}{la-onb-extension}
  $\R^n$ の正規直交系 $v_1,\dots,v_r$（$r\leq n$）は，
  正規直交基底 $v_1,\dots,v_n$ に延長できる．
  （$r=n$ のときは何もすることがない：正規直交系は一次独立
  （$\sum_ic_iv_i=0$ の両辺と $v_j$ の内積を取ると $c_j=0$）なので，
  $n$ 本あれば既に基底である．）
\end{claim}

\begin{proof}
  $r<n$ のとき $S:=\operatorname{span}\{v_1,\dots,v_r\}\neq\R^n$ だから
  $w\notin S$ が取れる．
  $w':=w-\sum_{i=1}^r\inner{v_i}{w}\,v_i$ とおくと，
  $w'=0$ なら $w\in S$ となり矛盾するから $w'\neq0$ であり，
  $\inner{v_i}{w'}=\inner{v_i}{w}-\inner{v_i}{w}=0$（$i\leq r$）．
  よって $v_{r+1}:=w'/\norm{w'}$ を加えれば正規直交系のまま
  1本延長できる．これを $n$ 本になるまで繰り返せばよい．
\end{proof}

\begin{theorem}{スペクトル分解}{la-spectral}
  対称行列 $A\in\R^{n\times n}$ に対し，直交行列 $U$ と
  実対角行列 $\Lambda=\diag(\lambda_1,\dots,\lambda_n)$ が存在して
  \[
    A=U\Lambda U^\top
  \]
  と書ける．$U$ の第 $i$ 列 $u_i$ は $Au_i=\lambda_iu_i$ をみたす
  固有ベクトルで，$u_1,\dots,u_n$ は $\R^n$ の正規直交基底をなす．
  外積の形では
  \[
    A=\sum_{i=1}^n\lambda_i\,u_iu_i^\top
  \]
  である．（線形代数の標準的な定理．証明は省略する．）
\end{theorem}

\begin{definition}{半正定値行列}{la-psd}
  対称行列 $A$ が\textbf{半正定値}であるとは，
  どの方向 $u\in\R^n$ でも二次形式が非負，すなわち
  \[
    u^\top Au\geq0\qquad(\forall u\in\R^n)
  \]
  が成り立つことをいい，$A\succeq0$ と書く．
  $u\neq0$ で常に $u^\top Au>0$ のとき\textbf{正定値}といい，$A\succ0$ と書く．
\end{definition}

\begin{claim}{半正定値性と固有値}{la-psd-eigen}
  対称行列 $A=\sum_i\lambda_iu_iu_i^\top$（スペクトル分解）について，
  $A\succeq0$ と「すべての固有値 $\lambda_i\geq0$」は同値である．
\end{claim}

\begin{proof}
  $A\succeq0$ ならば $\lambda_i=u_i^\top Au_i\geq0$．
  逆にすべての $\lambda_i\geq0$ ならば，任意の $u$ で
  $u^\top Au=\sum_i\lambda_i\inner{u_i}{u}^2\geq0$ である．
\end{proof}

\begin{claim}{合同変換は半正定値性を保つ}{la-congruence}
  $A\in\R^{n\times n}$ が対称で $A\succeq0$，$L\in\R^{k\times n}$ ならば，
  $LAL^\top\in\R^{k\times k}$ は対称かつ $LAL^\top\succeq0$ である．
\end{claim}

\begin{proof}
  $(LAL^\top)^\top=LA^\top L^\top=LAL^\top$ で対称．
  任意の $x\in\R^k$ で
  $x^\top(LAL^\top)x=(L^\top x)^\top A(L^\top x)\geq0$ である．
\end{proof}

\begin{claim}{固有値ごとにまとめたスペクトル分解}{la-spectral-grouped}
  対称行列 $M\in\R^{n\times n}$ について：
  \begin{enumerate}
    \item （存在）$M$ の相異なる固有値を $\nu_1,\dots,\nu_m$ とし，
      スペクトル分解 $M=\sum_i\lambda_iu_iu_i^\top$
      （定理~\ref{thm:la-spectral}）において
      $R_j:=\sum_{i:\lambda_i=\nu_j}u_iu_i^\top$ とおくと，
      各 $R_j$ は零でない対称行列で
      \[
        M=\sum_{j=1}^m\nu_jR_j,\qquad
        R_jR_l=\delta_{jl}R_j,\qquad
        \sum_{j=1}^mR_j=I.
      \]
    \item （一意性）逆に，相異なる実数 $\nu_1,\dots,\nu_m$ と
      零でない対称行列 $R_1,\dots,R_m$ が (i) の3条件をみたすならば，
      $\nu_1,\dots,\nu_m$ の全体は $M$ の固有値の全体と一致し，
      $\ker(M-\nu_jI)=\operatorname{ran}R_j
      :=\{R_jz\mid z\in\R^n\}$ であり，
      さらに各 $R_j$ は「$x$ をその $\ker(M-\nu_jI)$ 成分に写す行列」
      として $M$ だけから一意に定まる．
  \end{enumerate}
\end{claim}

\begin{proof}
  (i)：$R_j$ は対称行列の和ゆえ対称で，
  $\nu_j$ の固有ベクトル $u_i$ に対し $R_ju_i=u_i\neq0$ だから零でない．
  $\sum_jR_j=\sum_iu_iu_i^\top=I$
  （正規直交基底による展開 $x=\sum_iu_i\inner{u_i}{x}$），
  $M=\sum_j\nu_jR_j$ は分解のまとめ直しそのもの，そして正規直交性から
  \[
    R_jR_l
    =\sum_{i:\lambda_i=\nu_j}\;\sum_{i':\lambda_{i'}=\nu_l}
      u_i\inner{u_i}{u_{i'}}u_{i'}^\top
    =\delta_{jl}R_j.
  \]

  (ii)：まず $MR_j=\sum_l\nu_lR_lR_j=\nu_jR_j$ だから
  $\operatorname{ran}R_j\subset\ker(M-\nu_jI)$ である．
  逆に $Mx=\nu_jx$ とすると
  $0=(M-\nu_jI)x=\sum_l(\nu_l-\nu_j)R_lx$ であり，
  両辺に左から $R_{l'}$ を掛けると
  $(\nu_{l'}-\nu_j)R_{l'}x=0$，すなわち $l'\neq j$ で $R_{l'}x=0$．
  よって $x=\sum_lR_lx=R_jx\in\operatorname{ran}R_j$ となり，
  $\ker(M-\nu_jI)=\operatorname{ran}R_j$ が従う．
  同じ計算により，$Mx=\nu x$（$x\neq0$，$\nu\notin\{\nu_1,\dots,\nu_m\}$）
  ならすべての $l'$ で $R_{l'}x=0$，よって $x=\sum_lR_lx=0$ となり矛盾．
  各 $\nu_j$ が実際に固有値であることは $R_j\neq0$ と
  $\operatorname{ran}R_j\subset\ker(M-\nu_jI)$ から従うので，
  $M$ の固有値の全体は $\{\nu_1,\dots,\nu_m\}$ である．
  最後に，$x=\sum_ly_l$（$y_l\in\ker(M-\nu_lI)=\operatorname{ran}R_l$）と
  書けたとすると，$y_l=R_lz_l$ の形だから
  $R_jx=\sum_lR_jR_lz_l=R_jz_j=y_j$．
  すなわち $R_jx$ は $x$ の $\ker(M-\nu_jI)$ 成分として
  （$M$ の固有空間だけから）一意に決まるので，
  $R_j$ は $M$ だけから定まる．
\end{proof}

\section{半正定値平方根}

\begin{theorem}{半正定値平方根の存在と一意性}{la-sqrt}
  $\Sigma\succeq0$ に対し，
  $B^2=\Sigma$ をみたす対称行列 $B\succeq0$ がただ一つ存在する．
  これを $\Sigma^{1/2}$ と書く．
\end{theorem}

\begin{proof}
  \textbf{存在．}
  スペクトル分解（定理~\ref{thm:la-spectral}）で
  $\Sigma=U\Lambda U^\top$ と書くと，
  主張~\ref{clm:la-psd-eigen}よりすべての $\lambda_i\geq0$ である．
  $\Lambda^{1/2}:=\diag(\sqrt{\lambda_1},\dots,\sqrt{\lambda_n})$ として
  $B:=U\Lambda^{1/2}U^\top$ とおけば，$B$ は対称で，
  固有値 $\sqrt{\lambda_i}\geq0$ より $B\succeq0$
  （主張~\ref{clm:la-psd-eigen}），かつ
  $B^2=U\Lambda^{1/2}U^\top U\Lambda^{1/2}U^\top=U\Lambda U^\top=\Sigma$ である．

  \textbf{一意性．}
  $B\succeq0$ が対称で $B^2=\Sigma$ をみたすとする．
  主張~\ref{clm:la-spectral-grouped}(i) より，
  $B$ の相異なる固有値 $\mu_1,\dots,\mu_m$
  （$B\succeq0$ より $\mu_j\geq0$：主張~\ref{clm:la-psd-eigen}）と
  零でない対称行列 $Q_j$ で
  \[
    B=\sum_{j=1}^m\mu_jQ_j,
    \qquad
    Q_jQ_l=\delta_{jl}Q_j,
    \qquad
    \sum_{j=1}^mQ_j=I
  \]
  なるものが取れる．このとき
  \[
    \Sigma=B^2
    =\sum_{j,l}\mu_j\mu_l\,Q_jQ_l
    =\sum_{j=1}^m\mu_j^2\,Q_j
  \]
  であり，$\mu_j\geq0$ が相異なることから $\mu_j^2$ も相異なる．
  よって組 $\{(\mu_j^2,\,Q_j)\}_j$ は $\Sigma$ に対して
  主張~\ref{clm:la-spectral-grouped}(ii) の仮定をみたし，
  $\mu_j^2$（$=\Sigma$ の固有値の全体）と $Q_j$ は
  $\Sigma$ だけから一意に定まる．
  したがって $\mu_j=\sqrt{\mu_j^2}$（非負の平方根）も定まり，
  $B=\sum_j\mu_jQ_j$ は $\Sigma$ だけから一意に決まる．
\end{proof}

\begin{remark}{正定値行列の逆平方根}{la-inv-sqrt}
  $\Sigma\succ0$ のとき，固有値はすべて正
  （$\lambda_i=u_i^\top\Sigma u_i>0$）なので，
  $\Sigma^{1/2}=U\Lambda^{1/2}U^\top$ の固有値 $\sqrt{\lambda_i}$ も
  すべて正であり，$\Sigma^{1/2}$ は正則である．その逆行列を
  \[
    \Sigma^{-1/2}
    :=\bigl(\Sigma^{1/2}\bigr)^{-1}
    =U\Lambda^{-1/2}U^\top
    \qquad(\Lambda^{-1/2}:=\diag(1/\sqrt{\lambda_1},\dots,1/\sqrt{\lambda_n}))
  \]
  と書く（右の表示は $U\Lambda^{1/2}U^\top\cdot U\Lambda^{-1/2}U^\top=I$
  の直接計算による）．$\Sigma^{-1/2}$ は対称かつ正定値である．
\end{remark}

\section{作用素ノルムと Cauchy--Schwarz の不等式}

\begin{lemma}{半正定値形式の Cauchy--Schwarz}{la-psd-cs}
  $A\succeq0$ と $a,b\in\R^n$ に対し
  \[
    (a^\top Ab)^2\leq(a^\top Aa)(b^\top Ab).
  \]
  特に $A=I$ として，通常の Cauchy--Schwarz の不等式
  $\inner{a}{b}^2\leq\norm{a}^2\norm{b}^2$ を得る．
\end{lemma}

\begin{proof}
  $A$ は対称だから $b^\top Aa=a^\top Ab$ であり，任意の $t\in\R$ で
  \[
    0\leq(a+tb)^\top A(a+tb)
    =a^\top Aa+2t\,(a^\top Ab)+t^2\,(b^\top Ab).
  \]
  $b^\top Ab>0$ のとき：$t$ の2次式が常に非負だから判別式は非正，
  すなわち $(a^\top Ab)^2\leq(a^\top Aa)(b^\top Ab)$．
  $b^\top Ab=0$ のとき：$0\leq a^\top Aa+2t\,(a^\top Ab)$ が
  すべての $t$ で成り立つには $a^\top Ab=0$ が必要
  （さもなくば $t\to\pm\infty$ の一方で右辺が $-\infty$）．
  このとき両辺とも $0$ 以上・以下で不等式は成り立つ．
\end{proof}

\begin{claim}{Bessel の不等式}{la-bessel}
  $u_1,\dots,u_r$ を $\R^n$ の正規直交系とすると，任意の $w\in\R^n$ で
  \[
    \sum_{i=1}^r\inner{u_i}{w}^2\leq\norm{w}^2.
  \]
\end{claim}

\begin{proof}
  正規直交性から
  \[
    0\leq\Bigl\lVert w-\sum_{i=1}^r\inner{u_i}{w}\,u_i\Bigr\rVert^2
    =\norm{w}^2-\sum_{i=1}^r\inner{u_i}{w}^2.  \qedhere
  \]
\end{proof}

\begin{definition}{作用素ノルム}{la-opnorm}
  $C\in\R^{k\times n}$ の\textbf{作用素ノルム}を
  \[
    \norm{C}_{\mathrm{op}}
    :=
    \sup\{\norm{Cu}\mid u\in\R^n,\ \norm{u}\leq1\}
  \]
  と定める（$C$ が単位球をどれだけ引き伸ばすかの最大倍率）．
  成分ごとの Cauchy--Schwarz（補題~\ref{lem:la-psd-cs}）から
  $\norm{Cu}^2\leq\bigl(\sum_{i,j}c_{ij}^2\bigr)\norm{u}^2$ となるので，
  この上限は有限である．
  定義から $\norm{Cu}\leq\norm{C}_{\mathrm{op}}\norm{u}$（$\forall u$）であり，
  直交行列 $U$ はノルムを保つ（定義~\ref{def:la-symmetric}）から
  $\norm{U}_{\mathrm{op}}=1$ である．
\end{definition}

\section{特異値分解とトレース}

\begin{definition}{トレース}{la-trace}
  $A\in\R^{n\times n}$ の\textbf{トレース}を
  $\tr A:=\sum_{i=1}^na_{ii}$（対角成分の和）と定める．
\end{definition}

\begin{claim}{トレースの基本性質}{la-trace-basic}
  \begin{enumerate}
    \item $\tr$ は線形で，$\tr A^\top=\tr A$．
    \item （巡回性）$A\in\R^{k\times n}$，$B\in\R^{n\times k}$ に対し
      $\tr(AB)=\tr(BA)$．
    \item $u,v\in\R^n$ に対し $\tr(uv^\top)=\inner{u}{v}$．
    \item スペクトル分解 $A=U\Lambda U^\top$ に対し
      $\tr A=\sum_i\lambda_i$（固有値の和）．
  \end{enumerate}
\end{claim}

\begin{proof}
  (i) は定義から明らか．(ii) は
  \[
    \tr(AB)=\sum_{i=1}^k\sum_{j=1}^na_{ij}b_{ji}
    =\sum_{j=1}^n\sum_{i=1}^kb_{ji}a_{ij}=\tr(BA).
  \]
  (iii) は $\tr(uv^\top)=\sum_iu_iv_i=\inner{u}{v}$．
  (iv) は (ii) より
  $\tr(U\Lambda U^\top)=\tr(\Lambda U^\top U)=\tr\Lambda=\sum_i\lambda_i$．
\end{proof}

\begin{lemma}{特異値分解（SVD）}{la-svd}
  $A\in\R^{k\times n}$，$r:=\rank A$ とする．
  正数 $d_1\geq\cdots\geq d_r>0$ と，
  正規直交系 $u_1,\dots,u_r\in\R^k$，$v_1,\dots,v_r\in\R^n$ が存在して
  \[
    A=\sum_{i=1}^rd_i\,u_iv_i^\top
  \]
  と書ける．さらにこのとき
  \[
    AA^\top=\sum_{i=1}^rd_i^2\,u_iu_i^\top,
    \qquad
    A^\top A=\sum_{i=1}^rd_i^2\,v_iv_i^\top
  \]
  である．
\end{lemma}

\begin{proof}
  $AA^\top\in\R^{k\times k}$ は対称であり，
  $x^\top AA^\top x=\norm{A^\top x}^2\geq0$ より半正定値である．
  スペクトル分解（定理~\ref{thm:la-spectral}）で
  $AA^\top=\sum_{i=1}^k\lambda_iu_iu_i^\top$（$\{u_i\}$ は $\R^k$ の正規直交基底，
  $\lambda_1\geq\cdots\geq\lambda_k\geq0$：主張~\ref{clm:la-psd-eigen}）とし，
  $r':=\#\{i\mid\lambda_i>0\}$，$d_i:=\sqrt{\lambda_i}$（$i\leq r'$）とおく．

  $i\leq r'$ に対し $v_i:=A^\top u_i/d_i$ と定めると
  \[
    \inner{v_i}{v_j}
    =\frac{u_i^\top AA^\top u_j}{d_id_j}
    =\frac{\lambda_j\inner{u_i}{u_j}}{d_id_j}
    =\delta_{ij}
  \]
  だから $\{v_i\}_{i\leq r'}$ は正規直交系である．
  $A':=\sum_{i\leq r'}d_iu_iv_i^\top$ とおいて $A=A'$ を示す．
  正規直交基底 $\{u_i\}_{i=1}^k$ 上で転置を比較すると，
  $i\leq r'$ では
  \[
    A'^\top u_i=\sum_{j\leq r'}d_jv_j\inner{u_j}{u_i}=d_iv_i=A^\top u_i,
  \]
  $i>r'$ では $\norm{A^\top u_i}^2=u_i^\top AA^\top u_i=\lambda_i=0$ より
  $A^\top u_i=0=A'^\top u_i$．
  よって $A^\top=A'^\top$，すなわち $A=A'$ である．

  このとき $A=A'$ の表示から
  $Ax=\sum_{i\leq r'}d_i\inner{v_i}{x}\,u_i$ なので
  $\operatorname{ran}A\subset\operatorname{span}\{u_i\mid i\leq r'\}$，
  逆に $Av_i=d_iu_i$ より各 $u_i\in\operatorname{ran}A$（$i\leq r'$）．
  よって $\operatorname{ran}A=\operatorname{span}\{u_1,\dots,u_{r'}\}$ で，
  正規直交系の一次独立性（主張~\ref{clm:la-onb-extension}内の注意）から
  $\rank A=r'$，すなわち $r'=r$ である．
  最後に，$\{v_i\}$ の正規直交性から
  \[
    A^\top A
    =\Bigl(\sum_id_iv_iu_i^\top\Bigr)\Bigl(\sum_jd_ju_jv_j^\top\Bigr)
    =\sum_{i,j}d_id_jv_i\inner{u_i}{u_j}v_j^\top
    =\sum_{i=1}^rd_i^2\,v_iv_i^\top
  \]
  であり，$AA^\top=\sum_id_i^2u_iu_i^\top$ は構成そのものである．
\end{proof}

\begin{definition}{特異値}{la-singular}
  補題~\ref{lem:la-svd}の $d_1\geq\cdots\geq d_r>0$ に
  $0$ を並べて $\min\{k,n\}$ 個にしたものを $A$ の\textbf{特異値}とよび，
  $\sigma_1(A)\geq\cdots\geq\sigma_{\min\{k,n\}}(A)\geq0$ と書く．
  この定義は SVD の取り方によらない：
  補題~\ref{lem:la-svd}より $AA^\top=\sum_{i\leq r}d_i^2\,u_iu_i^\top$ で，
  $\{u_i\}$ を正規直交基底に延長（主張~\ref{clm:la-onb-extension}）すれば
  これは $AA^\top$ のスペクトル分解になるから，
  $d_i^2$ は重複度込みで $AA^\top$ の非零固有値と一致し
  （固有値が分解の係数に尽きることは
  主張~\ref{clm:la-spectral-grouped}(ii)），$A$ だけから定まる．
  同様に非零の特異値の2乗は $A^\top A$ の非零固有値とも一致するので，
  特に $\sigma_i(A)=\sigma_i(A^\top)$ である．
\end{definition}

\begin{lemma}{特異値の和とトレース}{la-trace-norm}
  $A\in\R^{k\times n}$ に対し
  \[
    \sum_i\sigma_i(A)
    =\tr\bigl((A^\top A)^{1/2}\bigr)
    =\tr\bigl((AA^\top)^{1/2}\bigr).
  \]
\end{lemma}

\begin{proof}
  SVD（補題~\ref{lem:la-svd}）の記号で，
  $B:=\sum_{i=1}^rd_i\,v_iv_i^\top$ とおく．
  $B$ は対称で，
  $x^\top Bx=\sum_id_i\inner{v_i}{x}^2\geq0$ より $B\succeq0$，
  かつ正規直交性より
  \[
    B^2=\sum_{i,j}d_id_j\,v_i\inner{v_i}{v_j}v_j^\top
    =\sum_{i=1}^rd_i^2\,v_iv_i^\top=A^\top A.
  \]
  平方根の一意性（定理~\ref{thm:la-sqrt}）より
  $(A^\top A)^{1/2}=B$ であり，
  主張~\ref{clm:la-trace-basic}(iii) から
  \[
    \tr\bigl((A^\top A)^{1/2}\bigr)
    =\sum_{i=1}^rd_i\,\tr(v_iv_i^\top)
    =\sum_{i=1}^rd_i\,\norm{v_i}^2
    =\sum_{i=1}^rd_i
    =\sum_i\sigma_i(A).
  \]
  $AA^\top$ についても同様である（$u_i$ を用いる）．
\end{proof}

\begin{lemma}{von Neumann 型トレース不等式}{la-vn-trace}
  $A,C\in\R^{n\times n}$，$\norm{C}_{\mathrm{op}}\leq1$ ならば
  \[
    \tr(CA)\leq\sum_i\sigma_i(A).
  \]
\end{lemma}

\begin{proof}
  SVD（補題~\ref{lem:la-svd}）で $A=\sum_{i=1}^rd_iu_iv_i^\top$ と書くと，
  トレースの線形性と巡回性（主張~\ref{clm:la-trace-basic}）から
  \[
    \tr(CA)
    =\sum_{i=1}^rd_i\,\tr(Cu_iv_i^\top)
    =\sum_{i=1}^rd_i\,\tr(v_i^\top Cu_i)
    =\sum_{i=1}^rd_i\,\inner{v_i}{Cu_i}.
  \]
  Cauchy--Schwarz（補題~\ref{lem:la-psd-cs}）と作用素ノルムの性質から
  \[
    \inner{v_i}{Cu_i}
    \leq\norm{v_i}\,\norm{Cu_i}
    \leq\norm{C}_{\mathrm{op}}\norm{u_i}
    \leq1
  \]
  だから，$\tr(CA)\leq\sum_id_i=\sum_i\sigma_i(A)$ である．
\end{proof}

\section{ブロック行列の半正定値性}

\begin{lemma}{ブロック半正定値行列の特徴づけ}{la-block-psd}
  $\Sigma_1,\Sigma_2\succeq0$（$n\times n$），$K\in\R^{n\times n}$ とする．
  次は同値である：
  \begin{enumerate}
    \item $2n\times2n$ の対称行列
      $M:=\begin{pmatrix}\Sigma_1&K\\K^\top&\Sigma_2\end{pmatrix}$
      が半正定値である．
    \item ある $C\in\R^{n\times n}$，$\norm{C}_{\mathrm{op}}\leq1$ が存在して
      $K=\Sigma_1^{1/2}C\,\Sigma_2^{1/2}$ と書ける．
  \end{enumerate}
  逆行列を使わないこの特徴づけにより，
  $\Sigma_1,\Sigma_2$ が特異（正則でない）場合も一様に扱える．
\end{lemma}

\begin{proof}
  $u,v\in\R^n$ に対し $w:=(u,v)\in\R^{2n}$ とおくと
  \[
    w^\top Mw
    =u^\top\Sigma_1u+2\,u^\top Kv+v^\top\Sigma_2v
  \]
  である（$v^\top K^\top u=u^\top Kv$ を用いた）．

  \textbf{(ii) $\Rightarrow$ (i)．}
  $a:=\Sigma_1^{1/2}u$，$b:=\Sigma_2^{1/2}v$ とおくと，
  $\Sigma_i^{1/2}$ の対称性から
  $u^\top Kv=u^\top\Sigma_1^{1/2}C\Sigma_2^{1/2}v=\inner{a}{Cb}$ であり，
  Cauchy--Schwarz（補題~\ref{lem:la-psd-cs}）と
  $\norm{Cb}\leq\norm{b}$ から
  \[
    w^\top Mw
    =\norm{a}^2+2\inner{a}{Cb}+\norm{b}^2
    \geq\norm{a}^2-2\norm{a}\,\norm{b}+\norm{b}^2
    =(\norm{a}-\norm{b})^2
    \geq0.
  \]

  \textbf{(i) $\Rightarrow$ (ii)．}
  $M\succeq0$ に Cauchy--Schwarz（補題~\ref{lem:la-psd-cs}）を
  ベクトル $(u,0)$，$(0,v)\in\R^{2n}$ に対して適用すると，
  $(u,0)^\top M(0,v)=u^\top Kv$，
  $(u,0)^\top M(u,0)=u^\top\Sigma_1u=\norm{\Sigma_1^{1/2}u}^2$，
  $(0,v)^\top M(0,v)=\norm{\Sigma_2^{1/2}v}^2$ だから
  \[
    (u^\top Kv)^2
    \leq
    \norm{\Sigma_1^{1/2}u}^2\,\norm{\Sigma_2^{1/2}v}^2
    \tag{$*$}
  \]
  を得る．

  $V_1:=\operatorname{ran}\Sigma_1^{1/2}:=\{\Sigma_1^{1/2}u\mid u\in\R^n\}$，
  $V_2:=\operatorname{ran}\Sigma_2^{1/2}$ とおく．
  スペクトル分解 $\Sigma_1^{1/2}=\sum_i\sqrt{\lambda_i}\,u_iu_i^\top$ から
  $V_1=\operatorname{span}\{u_i\mid\lambda_i>0\}$ である
  （$\subset$ は表示から，$\supset$ は $\lambda_i>0$ のとき
  $u_i=\Sigma_1^{1/2}(u_i/\sqrt{\lambda_i})$ から）．
  よって $\lambda_i>0$ に対応する $u_i$ たちが
  $V_1$ の正規直交基底 $e_1,\dots,e_m$ を与える．
  同様に $V_2$ の正規直交基底 $f_1,\dots,f_{m'}$ を取り，
  $V_2$ への\textbf{直交射影}を $P:=\sum_{j=1}^{m'}f_jf_j^\top$ と定める．
  定義から $Pw\in V_2$，$w\in V_2$ なら基底で展開して $Pw=w$，
  また Bessel の不等式（主張~\ref{clm:la-bessel}）より
  $\norm{Pw}^2=\bigl\lVert\sum_j\inner{f_j}{w}f_j\bigr\rVert^2
  =\sum_j\inner{f_j}{w}^2\leq\norm{w}^2$ である．
  なお $V_1=\{0\}$ または $V_2=\{0\}$ のときは，
  （$*$）ですべての $u,v$ について $u^\top Kv=0$，すなわち $K=O$ となり，
  $C=O$ が求める行列になるから，以下いずれも $\{0\}$ でないとする．

  $V_1\times V_2$ 上の関数 $\beta$ を
  \[
    \beta(\Sigma_1^{1/2}u,\;\Sigma_2^{1/2}v):=u^\top Kv
  \]
  で定める．これは well-defined である：
  $\Sigma_1^{1/2}u=\Sigma_1^{1/2}u'$ ならば（$*$）を $u-u'$ と $v$ に適用して
  $((u-u')^\top Kv)^2\leq\norm{\Sigma_1^{1/2}(u-u')}^2\norm{\Sigma_2^{1/2}v}^2=0$，
  すなわち $u^\top Kv=u'^\top Kv$ であり，$v$ 側も同様だからである．
  $\beta$ は双線形（$u^\top Kv$ が $u,v$ について線形で，
  値が代表 $u,v$ の取り方によらないことから）で，（$*$）より
  $\abs{\beta(a,b)}\leq\norm{a}\,\norm{b}$（$a\in V_1$，$b\in V_2$）をみたす．

  $b\in V_2$ に対し
  \[
    c(b):=\sum_{l=1}^m\beta(e_l,b)\,e_l\in V_1
  \]
  とおくと，$c$ は $b$ について線形で，$a\in V_1$ に対し
  $\inner{a}{c(b)}=\beta(a,b)$ をみたす（$a$ を基底で展開すればよい）．
  また $c(b)=0$ なら $\norm{c(b)}\leq\norm{b}$ は自明であり，
  $c(b)\neq0$ なら
  \[
    \norm{c(b)}^2=\inner{c(b)}{c(b)}=\beta(c(b),b)
    \leq\norm{c(b)}\,\norm{b}
  \]
  より $\norm{c(b)}\leq\norm{b}$ である．
  $C:=c\circ P$（$w\mapsto c(Pw)$；$c$ と $P$ は線形なので行列で表せる）
  とおけば，
  $\norm{Cw}=\norm{c(Pw)}\leq\norm{Pw}\leq\norm{w}$ より
  $\norm{C}_{\mathrm{op}}\leq1$ である．

  最後に $K=\Sigma_1^{1/2}C\Sigma_2^{1/2}$ を確かめる．
  任意の $u,v$ に対し，$\Sigma_2^{1/2}v\in V_2$ だから
  $P(\Sigma_2^{1/2}v)=\Sigma_2^{1/2}v$，すなわち
  $C(\Sigma_2^{1/2}v)=c(\Sigma_2^{1/2}v)$ であり，
  $\Sigma_1^{1/2}u\in V_1$ と合わせて
  \[
    u^\top\bigl(\Sigma_1^{1/2}C\Sigma_2^{1/2}\bigr)v
    =\inner{\Sigma_1^{1/2}u}{\,c(\Sigma_2^{1/2}v)}
    =\beta(\Sigma_1^{1/2}u,\;\Sigma_2^{1/2}v)
    =u^\top Kv
  \]
  であり，$u,v$ は任意だから両行列は一致する．
\end{proof}

\section{可換な半正定値行列}

\begin{lemma}{平方根は元の行列の多項式}{la-sqrt-poly}
  $\Sigma\succeq0$ に対し，実係数多項式 $p$ で
  $\Sigma^{1/2}=p(\Sigma)$ となるものが存在する．
\end{lemma}

\begin{proof}
  主張~\ref{clm:la-spectral-grouped}(i) より
  $\Sigma=\sum_j\mu_jQ_j$
  （$\mu_1,\dots,\mu_m$ は相異なる固有値で
  主張~\ref{clm:la-psd-eigen}より $\geq0$，
  $Q_jQ_l=\delta_{jl}Q_j$，$\sum_jQ_j=I$）と書ける．
  Lagrange 補間で $p(\mu_j)=\sqrt{\mu_j}$（$j=1,\dots,m$）をみたす
  多項式 $p$ を取る．
  $\Sigma^k=\sum_j\mu_j^kQ_j$（$k\geq1$）と $\sum_jQ_j=I$ から，
  多項式の線形結合として
  \[
    p(\Sigma)=\sum_jp(\mu_j)\,Q_j=\sum_j\sqrt{\mu_j}\,Q_j
  \]
  である．右辺は対称で，固有値 $\sqrt{\mu_j}\geq0$ より半正定値，
  さらに2乗すると $\sum_j\mu_jQ_j=\Sigma$ になるから，
  平方根の一意性（定理~\ref{thm:la-sqrt}）より
  $p(\Sigma)=\Sigma^{1/2}$ である．
\end{proof}

\begin{corollary}{可換な半正定値行列の平方根}{la-commuting-sqrt}
  $\Sigma_1,\Sigma_2\succeq0$ が $\Sigma_1\Sigma_2=\Sigma_2\Sigma_1$ をみたすならば：
  \begin{enumerate}
    \item $\Sigma_1^{1/2}$ と $\Sigma_2^{1/2}$ も可換である．
    \item $S:=\Sigma_1^{1/2}\Sigma_2^{1/2}$ は対称かつ半正定値である．
    \item $\bigl(\Sigma_1^{1/2}\Sigma_2\Sigma_1^{1/2}\bigr)^{1/2}
      =\Sigma_1^{1/2}\Sigma_2^{1/2}$．
  \end{enumerate}
\end{corollary}

\begin{proof}
  (i)：補題~\ref{lem:la-sqrt-poly}より
  $\Sigma_1^{1/2}=p(\Sigma_1)$，$\Sigma_2^{1/2}=q(\Sigma_2)$ と書ける．
  $\Sigma_1\Sigma_2=\Sigma_2\Sigma_1$ ならば
  $\Sigma_1$ と $\Sigma_2$ の任意の多項式どうしも可換だから，
  $\Sigma_1^{1/2}\Sigma_2^{1/2}=\Sigma_2^{1/2}\Sigma_1^{1/2}$ である．

  (ii)：対称性は (i) より
  $S^\top=\Sigma_2^{1/2}\Sigma_1^{1/2}=\Sigma_1^{1/2}\Sigma_2^{1/2}=S$．
  半正定値性：$\Sigma_2^{1/4}:=(\Sigma_2^{1/2})^{1/2}$ とおくと，
  補題~\ref{lem:la-sqrt-poly}を2回使って $\Sigma_2^{1/4}$ は
  $\Sigma_2$ の多項式だから $\Sigma_1^{1/2}$ と可換であり，
  \[
    S=\Sigma_1^{1/2}\Sigma_2^{1/4}\Sigma_2^{1/4}
    =\Sigma_2^{1/4}\,\Sigma_1^{1/2}\,\Sigma_2^{1/4}
    \succeq0
  \]
  （最後は合同変換：主張~\ref{clm:la-congruence}，
  $\Sigma_2^{1/4}$ は対称）．

  (iii)：(i) を繰り返し使うと
  \[
    S^2
    =\Sigma_1^{1/2}\Sigma_2^{1/2}\Sigma_1^{1/2}\Sigma_2^{1/2}
    =\Sigma_1^{1/2}\Sigma_2^{1/2}\Sigma_2^{1/2}\Sigma_1^{1/2}
    =\Sigma_1^{1/2}\Sigma_2\Sigma_1^{1/2}
  \]
  であり，(ii) と平方根の一意性（定理~\ref{thm:la-sqrt}）から
  $S=\bigl(\Sigma_1^{1/2}\Sigma_2\Sigma_1^{1/2}\bigr)^{1/2}$ である
  （左辺の行列が $\succeq0$ で平方根をもつことは
  主張~\ref{clm:la-congruence}による）．
\end{proof}
